\section{\texorpdfstring{\(p\)-Modular Systems}{p-Modular Systems}}

\begin{definition}[\(p\)-modular system]
    Let \(G\) be a finite group and let \(p\) be a prime number. A
    \textbf{\(p\)-modular system} for \(G\) is a triple
    \[
        (F,R,k)
    \]
    satisfying the following conditions:
    \begin{enumerate}
        \item \(F\) is a field of characteristic \(0\), equipped with a complete
        discrete valuation.
        \item \(R\) is the valuation ring of \(F\), with maximal ideal
        \[
            J(R)=(\pi).
        \]
        \item The residue field is
        \[
            k=R/(\pi),
        \]
        and
        \[
            \operatorname{char} k=p.
        \]
    \end{enumerate}
    If both \(F\) and \(k\) are splitting fields for \(G\), then \((F,R,k)\) is called
    a \textbf{splitting \(p\)-modular system} for \(G\).
\end{definition}

\begin{lemma}
    Let \(R\) be a complete discrete valuation ring with maximal ideal
    \[
        J(R)=(\pi),
    \]
    and residue field
    \[
        k=R/(\pi).
    \]
    Let \(G\) be a finite group. Then:
    \begin{enumerate}
        \item If \(S\) is a simple \(RG\)-module, then
        \[
            \pi S=0.
        \]

        \item The simple \(RG\)-modules are exactly the simple \(kG\)-modules,
        via the quotient homomorphism
        \[
            RG\longrightarrow kG.
        \]

        \item If \(U\) is an \(RG\)-module, then
        \[
            \pi U\subseteq \operatorname{Rad}(U).
        \]
        In particular,
        \[
            \pi RG\subseteq J(RG).
        \]

        \item For every \(RG\)-module \(U\),
        \[
            \operatorname{Rad}(U)/\pi U
            =
            \operatorname{Rad}(U/\pi U).
        \]
    \end{enumerate}
\end{lemma}

\begin{proof}
    First we prove that
    \[
        \pi RG\subseteq J(RG).
    \]
    Let \(z\in \pi RG\). Since \(RG\) is finite over the complete discrete valuation
    ring \(R\), it is complete with respect to the \(\pi\)-adic topology. Hence the
    geometric series
    \[
        1+z+z^2+\cdots
    \]
    converges in \(RG\), and gives an inverse of \(1-z\). Thus \(1-z\) is a unit.
    More generally, for every \(y\in RG\), one has \(yz\in \pi RG\), so
    \[
        1-yz
    \]
    is a unit. By the standard characterization of the Jacobson radical, this implies
    \[
        z\in J(RG).
    \]
    Hence
    \[
        \pi RG\subseteq J(RG).
    \]

    Now let \(S\) be a simple \(RG\)-module. Since the Jacobson radical annihilates
    every simple module, we have
    \[
        J(RG)S=0.
    \]
    Since \(\pi RG\subseteq J(RG)\), it follows that
    \[
        \pi S=0.
    \]
    This proves \((1)\).

    Since every simple \(RG\)-module is annihilated by \(\pi\), it factors through
    \[
        RG/\pi RG\cong kG.
    \]
    Therefore every simple \(RG\)-module is a simple \(kG\)-module. Conversely, every
    simple \(kG\)-module becomes a simple \(RG\)-module by restriction along the
    quotient map
    \[
        RG\longrightarrow kG.
    \]
    This proves \((2)\).

    Let \(U\) be an \(RG\)-module. If \(M\) is a maximal submodule of \(U\), then
    \(U/M\) is simple. By \((1)\),
    \[
        \pi(U/M)=0.
    \]
    Hence
    \[
        \pi U\subseteq M.
    \]
    Taking the intersection over all maximal submodules \(M\) of \(U\), we obtain
    \[
        \pi U\subseteq \operatorname{Rad}(U).
    \]
    Applying this to the regular module \(U=RG\), we get
    \[
        \pi RG\subseteq J(RG).
    \]
    This proves \((3)\).

    Finally, since every maximal submodule of \(U\) contains \(\pi U\), the maximal
    submodules of \(U/\pi U\) are precisely the quotients
    \[
        M/\pi U,
    \]
    where \(M\) runs through the maximal submodules of \(U\). Therefore
    \[
        \operatorname{Rad}(U/\pi U)
        =
        \bigcap_{M\text{ maximal in }U} M/\pi U
        =
        \left(\bigcap_{M\text{ maximal in }U}M\right)/\pi U
        =
        \operatorname{Rad}(U)/\pi U.
    \]
    This proves \((4)\).
\end{proof}

\begin{corollary}
    Let \(P\) and \(Q\) be finitely generated projective \(RG\)-modules. Then
    \[
        P\cong Q
        \quad\text{as }RG\text{-modules}
    \]
    if and only if
    \[
        P/\pi P\cong Q/\pi Q
        \quad\text{as }kG\text{-modules}.
    \]
\end{corollary}

\begin{proof}
    The forward implication is immediate by reducing modulo \(\pi\).

    Conversely, suppose that
    \[
        \overline{f}:P/\pi P\longrightarrow Q/\pi Q
    \]
    is an isomorphism of \(kG\)-modules. Let
    \[
        q_P:P\longrightarrow P/\pi P,
        \qquad
        q_Q:Q\longrightarrow Q/\pi Q
    \]
    be the natural quotient maps.

    Since \(P\) is projective and \(q_Q\) is surjective, the map
    \[
        \overline{f}\circ q_P:P\longrightarrow Q/\pi Q
    \]
    lifts to an \(RG\)-homomorphism
    \[
        f:P\longrightarrow Q
    \]
    such that
    \[
        q_Q\circ f=\overline{f}\circ q_P.
    \]
    Similarly, the inverse
    \[
        \overline{f}^{-1}:Q/\pi Q\longrightarrow P/\pi P
    \]
    lifts to an \(RG\)-homomorphism
    \[
        g:Q\longrightarrow P.
    \]

    Then the induced maps modulo \(\pi\) satisfy
    \[
        \overline{g f}
        =
        \overline{g}\,\overline{f}
        =
        \operatorname{id}_{P/\pi P},
    \]
    and similarly
    \[
        \overline{f g}
        =
        \operatorname{id}_{Q/\pi Q}.
    \]
    Hence
    \[
        (gf-\operatorname{id}_P)(P)\subseteq \pi P,
    \]
    and
    \[
        (fg-\operatorname{id}_Q)(Q)\subseteq \pi Q.
    \]

    Therefore the endomorphisms \(gf\) and \(fg\) are surjective by Nakayama's
    lemma, since
    \[
        \pi RG\subseteq J(RG).
    \]
    Since \(P\) and \(Q\) are finitely generated over the finite \(R\)-algebra \(RG\),
    they are finite \(R\)-modules. In particular, they are Noetherian. Hence a
    surjective endomorphism of \(P\) or \(Q\) is injective. Thus \(gf\) and \(fg\) are
    automorphisms.

    From \(gf\) being an automorphism, \(f\) is injective. From \(fg\) being an
    automorphism, \(f\) is surjective. Therefore \(f\) is an isomorphism:
    \[
        P\cong Q.
    \]
\end{proof}

